\section{Large Language Models}

\begin{frame}{Les Large Language Models (LLM)}

\textbf{Qu'est-ce qu'un LLM ?}
\begin{itemize}
    \item Non programmé avec des règles de grammaire
    \item Entraîné en lisant des milliards de textes (livres, sites web...)
    \item \textbf{Tâche principale :} \textcolor{red}{prédire} le mot suivant
    \item L'échelle compte : des milliards de paramètres
\end{itemize}

\vspace{0.5cm}

\onslide<2>
\textbf{De l'autocomplétion à la conversation :}
\begin{itemize}
    \item Modèle de base = autocomplétion performante
    \item Ajustement d'instructions = apprendre à suivre des demandes
    \item RLHF = apprendre à être utile grâce aux retours humains
\end{itemize}

\end{frame}

\begin{frame}{Ce que les LLM peuvent et ne peuvent pas faire}

\begin{columns}
\begin{column}{0.48\textwidth}
    \colormark{blue}\textbf{Capacités :}
    \begin{itemize}
        \item Reconnaître des patterns
        \item Générer du texte cohérent
        \item Raisonner sur du texte
        \item Résumer, traduire, reformuler
        \item Créer du contenu
    \end{itemize}
\end{column}

\begin{column}{0.48\textwidth}
    \onslide<2>
    \colorwarning{red}\textbf{Limites :}
    \begin{itemize}
        \item Pas de vraie compréhension
        \item Faits non fiables
        \item Pas d'info en temps réel
        \item Imprécisions
    \end{itemize}
\end{column}
\end{columns}

\vspace{0.5cm}

\textbf{Principe clé :} Toujours vérifier les informations critiques

\end{frame}

\begin{frame}{Pourquoi les LLM ont besoin d'aide}

\textbf{Trois problèmes majeurs :}

\vspace{0.3cm}

\begin{enumerate}
    \item<+-> \textbf{Ignorance des données spécifiques}
        \begin{itemize}
            \item Ne connaît pas les documents sur lesquels il n'a jamais vu
            \item Ne sait rien des spécificités liées à une entreprise particulière
        \end{itemize}
    
    \vspace{0.2cm}
    \item<+-> \textbf{Date limite de connaissances}
        \begin{itemize}
            \item Entraîné sur des données jusqu'à une certaine date
            \item Ignore les événements récents
        \end{itemize}
    
    \vspace{0.2cm}
    \item<+-> \textbf{Hallucinations}
        \begin{itemize}
            \item Invente des réponses plausibles quand il ne sait pas
            \item Présente la fiction comme un fait
        \end{itemize}
\end{enumerate}

\vspace{0.3cm}
\onslide<3>
\textbf{Solution :} Le RAG (Retrieval-Augmented Generation)

\end{frame}

\begin{frame}{RAG : Donner des documents au LLM}
\framesubtitle{Retrieval-Augmented Generation}

\textbf{Comment ça marche ?}

\vspace{0.3cm}

\begin{enumerate}
    \item[\textbf{1.}] L'utilisateur pose une question
    \item[\textbf{2.}] Le système cherche dans vos documents
    \item[\textbf{3.}] Les passages pertinents sont donnés au LLM
    \item[\textbf{4.}] Le LLM répond en se basant sur vos documents
\end{enumerate}

\vspace{0.5cm}

\textbf{Analogie :} Examen à livre ouvert

\vspace{0.5cm}
\onslide<2>
\textbf{Avantages :}
\begin{itemize}
    \item Réponses basées sur des données spécifiques
    \item Réduit les hallucinations
    \item Sources vérifiables
\end{itemize}

\end{frame}

\begin{frame}{Pourquoi le RAG est important pour les entreprises}

\textbf{Cas d'usage typiques :}

\vspace{0.3cm}

\begin{itemize}
    \item<+-> \textbf{Chatbots d'entreprise}
        \begin{itemize}
            \item Support client basé sur documentation produit
            \item Assistant RH pour les politiques internes
        \end{itemize}

    \item<+-> \textbf{Analyse documentaire à grande échelle}
        \begin{itemize}
            \item Recherche dans des milliers de contrats
            \item Synthèse de rapports techniques
        \end{itemize}

    \item<+-> \textbf{Conformité et réglementation}
        \begin{itemize}
            \item Q\&A sur les procédures de sécurité
            \item Vérification de conformité réglementaire
        \end{itemize}
\end{itemize}

\vspace{0.3cm}
\onslide<3>
\textbf{Bénéfice clé :} Démocratiser l'accès à l'information spécialisée

\end{frame}

\begin{frame}{Agents IA : au-delà de la conversation}

\textbf{Quelle différence ?}

\vspace{0.3cm}

\begin{tabular}{p{5cm}|p{5cm}}
\textbf{Chatbot classique} & \textbf{Agent IA} \\
\hline
Répond aux questions & Exécute des tâches \\
Conversation passive & Actions autonomes \\
Une interaction à la fois & Planification multi-étapes \\
Pas d'outils externes & Utilise calculatrice, web, APIs \\
\end{tabular}

\vspace{0.5cm}

\textbf{Exemple :}
\begin{itemize}
    \item Chatbot : "Le train Paris-Lyon coûte environ 100€"
    \item Agent : Cherche vols → Compare prix → Réserve → Envoie confirmation
\end{itemize}

\end{frame}

\begin{frame}{Composants d'un agent IA}

\textbf{5 capacités essentielles :}

\vspace{0.3cm}

\begin{enumerate}
    \item<+-> \textbf{Compréhension de l'objectif}
        \begin{itemize}
            \item Interpréter la demande de l'utilisateur
            \item Identifier le résultat souhaité
        \end{itemize}
    
    \item<+-> \textbf{Planification}
        \begin{itemize}
            \item Décomposer la tâche en étapes
            \item Décider de l'ordre des actions
        \end{itemize}
    
    \item<+-> \textbf{Utilisation d'outils}
        \begin{itemize}
            \item Calculatrice, recherche web, base de données
            \item APIs externes (météo, réservation, etc.)
        \end{itemize}
    
    \item<+-> \textbf{Mémoire}
        \begin{itemize}
            \item Garder le contexte de la conversation
            \item Apprendre des interactions précédentes
        \end{itemize}
    
    \item<+-> \textbf{Exécution et itération}
        \begin{itemize}
            \item Agir, observer le résultat, ajuster
        \end{itemize}
\end{enumerate}

\end{frame}

\begin{frame}{Défis des agents autonomes}

\textbf{Limitations actuelles :}

\vspace{0.3cm}

\begin{itemize}
    \item<+-> \textbf{Fiabilité}
        \begin{itemize}
            \item Erreurs possibles dans les tâches multi-étapes
            \item Peut mal interpréter les instructions
        \end{itemize}
    
    \item<+-> \textbf{Sécurité}
        \begin{itemize}
            \item Agents avec accès aux systèmes = risques
            \item Besoin de contrôles et limites strictes
        \end{itemize}
    
    \item<+-> \textbf{Coût}
        \begin{itemize}
            \item Nombreux appels au LLM = coûteux
            \item Temps d'exécution parfois long
        \end{itemize}
    
    \item<+-> \textbf{Transparence}
        \begin{itemize}
            \item Difficile de comprendre le raisonnement
            \item Audit et explicabilité limités
        \end{itemize}
\end{itemize}

\end{frame}

\begin{frame}{Synthèse : trois niveaux d'IA}

\begin{table}
\small
\begin{tabular}{l|p{3.5cm}|p{3.5cm}}
\textbf{Technologie} & \textbf{Fonction} & \textbf{Usage} \\
\hline
\textbf{LLM} & Comprendre et générer du texte & Rédaction, traduction, analyse \\
\hline
\textbf{RAG} & LLM + vos documents & Chatbot entreprise, recherche documentaire \\
\hline
\textbf{Agent} & RAG + actions autonomes & Automatisation de tâches complexes \\
\hline
\end{tabular}
\end{table}

\vspace{0.5cm}

\textbf{Progression naturelle :}
\begin{center}
LLM → RAG → Agent
\end{center}

\vspace{0.3cm}

Chaque niveau ajoute des capacités mais aussi des risques à gérer

\end{frame}